\begin{lemma}
Let $G$ be a finite $\Delta$-regular graph and let $0<\gamma\le\Delta$.
Then there is a probability mass function $\mu$ on subsets of $V(G)$ which is
the activation--priority independent-set sampling law of
\noderef{RandomIndependentSetSampling} with activation probability
$\gamma/\Delta$.
\end{lemma}
\begin{proof}
Put $p=\gamma/\Delta$.  The assumptions give $0<p\le 1$.  On the finite
space of activation sets, give each set $A\subseteq V(G)$ the Bernoulli
product weight
\[
p^{|A|}(1-p)^{|V(G)|-|A|}.
\]
These weights are nonnegative and sum to $(p+(1-p))^{|V(G)|}=1$ by
\noderef{SamplingBernoulliActivationMass}, the finite binomial/product
expansion over the independent activation decisions.
Independently of this activation choice, take the priority vector
\(\pi\in[0,1]^{V(G)}\) with the product Lebesgue probability measure.  Because
\(V(G)\) is finite, this product measure exists on the finite Cartesian
product of unit intervals; each coordinate is uniform on $[0,1]$, and for
every vector \(t\in[0,1]^{V(G)}\) the lower rectangle has measure
\(\prod_v t_v\).  The existence of the joint activation--priority law
\(\nu\), with the Bernoulli atom formula, unit-cube support on each activation
atom, lower-rectangle product identity, and fixed-activation comparison
identity, is exactly \noderef{SamplingFiniteProductPriorityLawExists}.

Let \(\nu\) be the law supplied by
\noderef{SamplingFiniteProductPriorityLawExists}.  Then \(\nu\) has total mass
one.  Its activation atoms have exactly the displayed Bernoulli weights,
priorities lie in the unit cube with full conditional measure on every
activation atom, and the lower-rectangle identity on each atom is the product
formula for independent uniform coordinates.
For a fixed activation set \(A\) and a fixed activated vertex \(v\), the event
that \(v\) beats all activated neighbors is computed by conditioning on
\(\pi(v)=a\): each activated neighbor has priority below \(a\) with
probability \(a\), independently.  Integrating over \(a\in[0,1]\) gives
\[
\Pr(A\text{ is the activation set and }v\text{ beats its activated
neighbors})
=\Pr(A)\int_0^1 a^{|A\cap N_G(v)|}\,da,
\]
which is precisely the fixed-activation comparison clause in
\noderef{RandomIndependentSetSampling}.

Define
\[
I(A,\pi)=\{v\in A:\pi(v)>\pi(w)\text{ for every }w\in A\cap N_G(v)\}
\]
and apply \noderef{SamplingActivationPriorityPushForwardSupport}.  This
supplies a finite mass function \(\mu\) on subsets, proves that \(\mu\) is
nonnegative and has total mass one, identifies \(\mu(S)\) with
\(\nu\{(A,\pi):I(A,\pi)=S\}\), and proves that every set with positive
\(\mu\)-mass is independent.  The last assertion follows from the strict
priority rule: if adjacent vertices \(u\) and \(v\) were both in
\(I(A,\pi)\), then the two inequalities \(\pi(v)<\pi(u)\) and
\(\pi(u)<\pi(v)\) would both have to hold.

It remains only to check the two-vertex integral clause of
\noderef{RandomIndependentSetSampling}.  Let \(u\ne v\) be nonadjacent.  By
\noderef{SamplingNonadjacentPairChamberIntegral}, applied to the same
\(\nu\), the push-forward \(\mu\), the product clauses above, and the
\(\Delta\)-regularity of \(G\), the probability that the output contains both
\(u\) and \(v\) is exactly the displayed chamber integral.  Concretely, for
both \(u\) and \(v\) to survive, they must both be activated, contributing
the factor \(p^2\).  Write the rescaled priority coordinates as
\(x=\gamma(1-\pi(u))\) and \(y=\gamma(1-\pi(v))\).  Splitting the square
\([0,\gamma]^2\) into the two ordered chambers and using symmetry gives the
factor \(2/\Delta^2\) and the region \(0\le x\le y\le\gamma\).  On this
chamber, a neighbor of \(u\) which is not also a neighbor of \(v\) avoids
blocking \(u\) with probability \(1-x/\Delta\), while every neighbor of
\(v\) avoids blocking \(v\) with probability \(1-y/\Delta\).  Common
neighbors are counted in the second group only on this ordered chamber, so
the first exponent is
\(\Delta-|N_G(u)\cap N_G(v)|\) and the second exponent is \(\Delta\).  The
regularity of \(G\) supplies the two degree counts equal to \(\Delta\).
Multiplying the independent avoidance probabilities and integrating over the
ordered chamber gives exactly
\[
\frac{2}{\Delta^2}\int_0^\gamma\int_x^\gamma
\left(1-\frac{x}{\Delta}\right)^{
\Delta-|N_G(u)\cap N_G(v)|}
\left(1-\frac{y}{\Delta}\right)^\Delta\,dy\,dx .
\]
Therefore the measure \(\nu\) and its push-forward \(\mu\) satisfy every
clause in \noderef{RandomIndependentSetSampling}, proving the claimed
existence.
\end{proof}
